\Article{4}
Daraus lässt sich leicht ersehen, welche magnetische Kraft die Erde in jedem Punkt des Raums ausübt, wie auch immer die magnetischen Flüssigkeiten in ihr verteilt sein mögen. Man denke sich das gesamte Erdvolumen, soweit es freien Magnetismus, das heißt getrennte magnetische Flüssigkeiten, enthält, in unendlich kleine Elemente zerlegt. Mit $\mathrm{d}\mu$ bezeichne man allgemein die in jedem Element enthaltene Menge freier magnetischer Flüssigkeit, wobei südliche Flüssigkeit stets als negativ gilt. Ferner sei $\rho$ die Entfernung des $\mathrm{d}\mu$ von einem beliebigen Punkt des Raums mit den rechtwinkligen Koordinaten $x,y,z$. Schließlich sei $V$ die Summe der Terme $\frac{\mathrm{d}\mu}{\rho}$ mit umgekehrtem Vorzeichen, erstreckt über alle magnetischen Teilchen der Erde. oder es sei
\SourceLeaf{132}
\[
V=-\int\frac{\mathrm{d}\mu}{\rho}
\]
Damit hat $V$ in jedem Punkt des Raums einen bestimmten Wert. Es ist also eine Funktion von $x,y,z$ oder von beliebigen drei anderen veränderlichen Größen, durch die sich die Punkte des Raums unterscheiden lassen. Die magnetische Kraft $\psi$ in jedem Punkt des Raums und die Komponenten $\xi,\eta,\zeta$, die sich aus der Zerlegung von $\psi$ parallel zu den Koordinatenachsen ergeben, erhält man dann aus den Formeln
\[
\xi=\frac{\mathrm{d}V}{\mathrm{d}x},\quad \eta=\frac{\mathrm{d}V}{\mathrm{d}y},\quad \zeta=\frac{\mathrm{d}V}{\mathrm{d}z},\quad \psi=\BareRoot{(\xi\xi+\eta\eta+\zeta\zeta)}
\]

\Article{5}
Zunächst sollen einige allgemeine Sätze hergeleitet werden, die von der Form der Funktion $V$ unabhängig und wegen ihrer Einfachheit und Eleganz bemerkenswert sind.

Das vollständige Differential von $V$ lautet
\[
\begin{aligned}
\mathrm{d}V&=\frac{\mathrm{d}V}{\mathrm{d}x}\cdot\mathrm{d}x+\frac{\mathrm{d}V}{\mathrm{d}y}\cdot\mathrm{d}y+\frac{\mathrm{d}V}{\mathrm{d}z}\cdot\mathrm{d}z\\
&=\xi\,\mathrm{d}x+\eta\,\mathrm{d}y+\zeta\,\mathrm{d}z
\end{aligned}
\]

Bezeichnet man mit $\mathrm{d}s$ den Abstand zwischen den beiden Punkten, auf die sich $V$ und $V+\mathrm{d}V$ beziehen, und mit $\theta$ den Winkel, den die Richtung der magnetischen Kraft $\psi$ mit $\mathrm{d}s$ bildet, so erhält man
\[
\mathrm{d}V=\psi\cos\theta\cdot\mathrm{d}s
\]
denn $\frac{\xi}{\psi},\frac{\eta}{\psi},\frac{\zeta}{\psi}$ sind die Kosinus der Winkel, die die Richtung von $\psi$ mit den Koordinatenachsen bildet, während $\frac{\mathrm{d}x}{\mathrm{d}s},\frac{\mathrm{d}y}{\mathrm{d}s},\frac{\mathrm{d}z}{\mathrm{d}s}$ die Kosinus der Winkel zwischen $\mathrm{d}s$ und denselben Achsen sind. Somit ist $\frac{\mathrm{d}V}{\mathrm{d}s}$ gleich der auf die Richtung von $\mathrm{d}s$ projizierten Kraft. Dasselbe folgt bereits aus der Gleichung $\frac{\mathrm{d}V}{\mathrm{d}x}=\xi$, wenn man bedenkt, dass die Koordinatenachsen beliebig gewählt werden können.

\Article{6}
Verbindet man zwei Punkte im Raum, $P^0,P'$, durch eine beliebige Linie, von der $\mathrm{d}s$ ein beliebiges Element darstellt, und bezeichnet wie zuvor $\theta$ den Winkel zwischen $\mathrm{d}s$ und der Richtung der dort wirkenden \SourceLeaf{133}magnetischen Kraft und $\psi$ deren Intensität, so gilt
\[
\int\psi\cos\theta\cdot\mathrm{d}s=V'-V^0
\]
wenn man die Integration über die gesamte Linie erstreckt und mit $V^0,V'$ die Werte von $V$ an den Endpunkten bezeichnet.

Die folgenden Folgerungen aus diesem fruchtbaren Satz verdienen besondere Erwähnung.

I. Das Integral $\int\psi\cos\theta\cdot\mathrm{d}s$ hat auf jedem Weg von $P^0$ nach $P'$ denselben Wert.

II. Das Integral $\int\psi\cos\theta\cdot\mathrm{d}s$ über die gesamte Länge einer beliebigen in sich zurückkehrenden Linie ist stets $=0$.

III. Auf einer geschlossenen Linie muss, sofern nicht überall $\theta=90^0$ gilt, ein Teil der Werte von $\theta$ kleiner und ein Teil größer als $90^0$ sein.

\Article{7}
\expandafter\def\csname HistoricalNoteText15\endcsname{Könnte die Funktion $V$ jede beliebig gewählte Form haben, so könnte in besonderen Fällen ein Maximum oder Minimum von $V$ einem isolierten Punkt oder einer isolierten Linie entsprechen, wobei ringsum nur kleinere oder nur größere Werte vorkämen. Es könnte auch einer Fläche entsprechen, auf deren \textit{beiden} Seiten zugleich kleinere oder größere Werte vorkämen. Die Bedingungen, denen die Funktion $V$ unterliegt, lassen diese Ausnahmefälle jedoch nicht zu. Eine ausführliche Behandlung dieses Gegenstands ist für unseren gegenwärtigen Zweck entbehrlich und muss einer anderen Gelegenheit vorbehalten bleiben.}
Die Fläche, auf der $V$ in jedem Punkt denselben bestimmten Wert $=V^0$ hat, trennt die Punkte des Raums, in denen $V$ größer als $V^0$ ist, von denen, in denen es kleiner als $V^0$ ist\HistoricalNote{15}. Aus dem Satz in Artikel 5 folgt leicht, dass die magnetische Kraft in jedem Punkt dieser Fläche senkrecht zur Fläche gerichtet ist, und zwar zu der Seite hin, auf der die größeren Werte von $V$ auftreten. Ist $\mathrm{d}s$ eine unendlich kleine, zur Fläche senkrechte Linie und $V^0+\mathrm{d}V^0$ der Wert von $V$ an ihrem anderen Endpunkt, so ist die Intensität der magnetischen Kraft $=\frac{\mathrm{d}V^0}{\mathrm{d}s}$. Die Gesamtheit der Punkte, für die $V=V^0+\mathrm{d}V^0$ gilt, bildet eine zweite, der ersten unendlich nahe Fläche. An den verschiedenen Stellen des gesamten Zwischenraums ist die Intensität der magnetischen Kraft umgekehrt proportional zum Abstand der beiden Flächen. Lässt man $V$ in unendlich kleinen, aber gleichen Schritten variieren, so entsteht ein System \SourceLeaf{134}von Flächen, die den Raum in unendlich dünne Schichten teilen. Die Dicke der Schichten ist dann sowohl an verschiedenen Stellen ein und derselben Schicht als auch in verschiedenen Schichten umgekehrt proportional zur Intensität der magnetischen Kraft.

\Article{8}
Wir wollen nun betrachten, wie sich die Werte von $V$ auf der Erdoberfläche verhalten.

In einem Punkt $P$ der Erdoberfläche seien $\psi$ die Intensität und $PM$ die Richtung der gesamten magnetischen Kraft. Ferner seien $\omega$ die Intensität und $PN$ die Richtung der auf die Horizontalebene projizierten Kraft. Gleichbedeutend damit ist $PN$ die Richtung des magnetischen Meridians, und zwar vom Südpol der Magnetnadel zum Nordpol. Weiter sei $i$ der Winkel zwischen $PM$ und $PN$, also die Inklination, und $\theta,t$ seien die Winkel zwischen dem Element $\mathrm{d}s$ einer auf der Erdoberfläche liegenden Linie und den Richtungen $PM,PN$. Schließlich entsprechen $V$ und $V+\mathrm{d}V$ dem Anfangs- und dem Endpunkt von $\mathrm{d}s$. Folglich gilt
\[
\cos\theta=\cos i\cos t,\qquad \omega=\psi\cos i
\]
und die Gleichung aus Artikel 5 wird zu
\[
\mathrm{d}V=\omega\cos t\cdot\mathrm{d}s
\]

Sind also zwei Punkte $P^0,P'$ auf der Erdoberfläche, in denen $V$ die Werte $V^0,V'$ hat, durch eine ganz auf der Erdoberfläche liegende Linie verbunden, von der $\mathrm{d}s$ ein beliebiges Element bezeichnet, so gilt
\[
\int\omega\cos t\cdot\mathrm{d}s=V'-V^0
\]
wenn die Integration über die gesamte Linie erstreckt wird. Offenbar gelten auch hier drei ganz ähnliche Folgerungen wie in Artikel 6.

I. Das Integral $\int\omega\cos t\,\mathrm{d}s$ hat denselben Wert, auf welchem Weg man auch auf der Erdoberfläche von $P^0$ nach $P'$ gelangt.

II. Das Integral $\int\omega\cos t\cdot\mathrm{d}s$ über die gesamte Länge einer auf der Erdoberfläche liegenden geschlossenen Linie ist stets $=0$.

III. Auf einer solchen geschlossenen Linie muss, falls nicht \SourceLeaf{135}überall $t=90^0$ gilt, ein Teil der Werte von $t$ spitz und ein Teil stumpf sein.

\Article{9}
Die Sätze I und II des vorigen Artikels (eigentlich nur zwei verschiedene Fassungen derselben Sache) lassen sich zumindest näherungsweise anhand tatsächlicher Beobachtungen prüfen. Es sei $P^0P'P''....P^0$ ein Polygon auf der Erdoberfläche, dessen Seiten die kürzesten Linien zwischen ihren Endpunkten sind. Betrachtet man die Erde hier als Kugel, so sind es also Großkreisbögen. Es seien $\omega^0,\omega',\omega''$ usw. die Intensitäten der horizontalen magnetischen Kraft in den Punkten $P^0,P',P''$ usw. Weiter seien $\delta^0,\delta',\delta''$ usw. die Deklinationen, die man wie üblich westlich von Nord als positiv und östlich davon als negativ ansetzen kann. Schließlich sei $(01)$ das Azimut der Linie $P^0P'$ in $P^0$, wie üblich von Süden über Westen gezählt. Ebenso sei $(10)$ das Azimut derselben Linie in umgekehrter Richtung in $P'$ usw.

Man beachte. dass sich $t$ zwar auf jeder Seite des Polygons stetig, an den Eckpunkten dagegen sprunghaft ändert und dort daher zwei verschiedene Werte hat. In $P'$ hat $t$ zum Beispiel

den Wert $(10)+\delta'$, wenn $P'$ der Endpunkt von $P^0P'$ ist,

den Wert $180^0+(12)+\delta'$, wenn $P'$ der Anfangspunkt von $P'P''$ ist.

Als Näherungswert des Integrals $\int\omega\cos t\cdot\mathrm{d}s$ über $P^0P'$ kann man
\[
\tfrac12(\omega^0\cos t^0+\omega'\cos t')\cdot P^0P'
\]
nehmen, wenn $t^0,t'$ die Werte von $t$ in $P^0$ als Anfangspunkt und in $P'$ als Endpunkt von $P^0P'$ bezeichnen. Mehr als diese Näherung lässt sich nicht gewinnen, solange die Werte von $\omega$ und $t$ nur an den Endpunkten $P^0,P'$ vorliegen. Sie ist umso zulässiger, je kürzer die Linie ist. Der angegebene Ausdruck lautet in unserer Bezeichnung
\[
=\tfrac12\{\omega'\cos((10)+\delta')-\omega^0\cos((01)+\delta^0)\}\cdot P^0P'
\]
Entsprechend ist der Näherungswert des Integrals über $P'P''$
\SourceLeaf{136}
\[
=\tfrac12\{\omega''\cos((21)+\delta'')-\omega'\cos((12)+\delta')\}\cdot P'P''
\]
und so fort über das ganze Polygon.

Für ein Dreieck ergibt unser Satz also die näherungsweise richtige Gleichung
\[
\begin{aligned}
&\omega^0\{P^0P'\cos((01)+\delta^0)-P^0P''\cos((02)+\delta^0)\}\\
{}+{}&\omega'\{P'P''\cos((12)+\delta')-P^0P'\cos((10)+\delta')\}\\
{}+{}&\omega''\{P^0P''\cos((20)+\delta'')-P'P''\cos((21)+\delta'')\}\\
={}&0
\end{aligned}
\]

Offenbar sind die Einheiten für die Intensitäten und Abstände in dieser Gleichung beliebig.

